(a) Equivalently, $\overline I=\bigcap_{i=0}^r(IS_{x_i}\cap S)$. Each contraction is an ideal, so $\overline I$ is an ideal. If $s\in\overline I$, a sufficiently large power of each $x_i$ carries $s$ into $I$. Since $I$ is homogeneous, the same is true of every homogeneous component of $s$. Thus $\overline I$ is homogeneous.

(b) On $D_+(x_i)$, the ideal sheaf defined by $I$ is $I_{(x_i)}$. For homogeneous $s\in S_d$, one has $s/x_i^d\in I_{(x_i)}$ exactly when $x_i^ns\in I$ for some $n$. Consequently the homogeneous elements of $\overline I$ are precisely those whose associated sections belong locally to this ideal sheaf on every $D_+(x_i)$. Equality of ideal sheaves therefore implies equality of saturations. Conversely, $I$ and $\overline I$ have identical localized ideals, so equal saturations give the same subscheme.

(c) Let $J=\bigoplus_{d\ge0}\Gamma(X,\mathcal I_Y(d))\subseteq S$. A homogeneous $s\in S_d$ belongs to $\overline J$ exactly when $s/x_i^d$ belongs to $\mathcal I_Y(D_+(x_i))$ for every $i$. This is exactly the condition that the section $s$ of $\mathcal O_X(d)$ belongs to $\Gamma(X,\mathcal I_Y(d))$. Thus $J=\overline J$. Every other defining ideal has the same saturation by (b), hence is contained in $J$.

(d) Every closed subscheme arises from a homogeneous ideal, and (b) identifies precisely the ideals giving the same subscheme. Part (c) gives its unique saturated representative.
